\chapter{Hip Dysplasia}
\index{Hip Dysplasia}

\section*{Definition and Overview}
Developmental dysplasia of the hip (DDH), commonly called hip dysplasia, is a condition in which the hip joint does not develop normally, so the ball of the thigh bone (femoral head) does not sit firmly in the hip socket (acetabulum). The socket may be too shallow, and in severe cases the hip can dislocate. Hip dysplasia is present from birth or early life, is more common in girls, and is more likely in babies born breech or with a family history. Early detection and treatment produce excellent outcomes, while untreated dysplasia causes pain, limping, and early arthritis.

\section*{Epidemiology}
Hip dysplasia affects approximately 1--3 babies in every 1,000, with the risk rising to around 3--5\% in girls born breech. The left hip is more commonly affected than the right, and both hips are involved in about a fifth of cases. In many countries, all newborns are screened by clinical examination, and at-risk babies have ultrasound imaging. Late diagnosis is uncommon now but still occurs, and untreated adult dysplasia is a significant cause of early hip arthritis and hip replacement before age 60.

\section*{Aetiology and Pathophysiology}
\begin{itemize}
    \item \textbf{Congenital factors:} laxity of the ligaments and hormones in pregnancy, more common in first-born children
    \item \textbf{Breech position:} strong risk factor, especially in girls
    \item \textbf{Family history:} the condition runs in families
    \item \textbf{Oligohydramnios:} low amniotic fluid restricts fetal movement
    \item \textbf{Pathophysiology:} the femoral head is unstable in the shallow or displaced socket; with weight-bearing, the socket fails to deepen normally, leading to instability, dislocation, and abnormal joint loading
    \item \textbf{Sequelae:} if untreated, the hip develops arthritis early in adult life
\end{itemize}

\section*{Clinical Features}
\begin{itemize}
    \item Limited abduction of the affected hip in a baby
    \item Asymmetry of thigh or buttock creases
    \item A leg that appears shorter on the affected side
    \item A "clunk" felt when the hip is moved during examination
    \item In walking-age children: a limp or waddling gait
    \item In adults: groin or hip pain, limping, and reduced range of movement
    \item Many babies show no obvious signs, which is why screening matters
\end{itemize}

\section*{Differential Diagnosis}
\begin{itemize}
    \item Other causes of hip instability and limb asymmetry in babies
    \item Septic arthritis of the hip, which causes fever and a painful, irritable hip
    \item Legg--Calvé--Perthes disease in older children
    \item Slipped capital femoral epiphysis in adolescents
    \item Neuromuscular conditions causing abnormal gait
    \item Ultrasound and clinical assessment distinguish these conditions
\end{itemize}

\section*{When to Seek Medical Care}
All babies are screened at birth and at routine health checks, and parents should report any concerns about leg length, hip movement, or gait. Older children and adults with hip or groin pain should be assessed, especially with a history of risk factors.

\textbf{Contact your local emergency services for:}
\begin{itemize}
    \item A baby who refuses to move a leg or is distressed with nappy changes
    \item Fever with a painful, swollen hip (possible infection)
    \item Inability to bear weight after a fall
    \item Sudden severe hip or leg pain
\end{itemize}

\section*{Diagnosis and Investigations}
\begin{itemize}
    \item \textbf{Clinical examination:} all newborns are examined for hip stability and range of movement
    \item \textbf{Ultrasound:} the standard imaging test in babies under 6 months, showing socket depth and hip position
    \item \textbf{X-ray:} used after 4--6 months when the hip bones have begun to ossify
    \item \textbf{Assessment of at-risk babies:} imaging for breech birth, family history, or abnormal examination
    \item \textbf{Adult imaging:} X-ray and MRI for pain and arthritis assessment
\end{itemize}

\section*{Management}
\begin{itemize}
    \item \textbf{Observation:} mild, stable dysplasia may be monitored with repeat ultrasound
    \item \textbf{Harness treatment:} the Pavlik harness holds the hips in the correct position for 6--12 weeks in young babies
    \item \textbf{Closed reduction:} gentle repositioning of the hip with a cast or brace
    \item \textbf{Surgery:} open reduction and pelvic osteotomy for older children or failed conservative treatment
    \item \textbf{Adult management:} physiotherapy, activity modification, and pain management for early arthritis
    \item \textbf{Hip preservation surgery:} periacetabular osteotomy in young adults with dysplasia
    \item \textbf{Hip replacement:} for advanced arthritis
\end{itemize}

\section*{Complications}
\begin{itemize}
    \item Persistent instability or redislocation after treatment
    \item Avascular necrosis of the femoral head from treatment or persistent dislocation
    \item Early hip arthritis in adulthood
    \item Leg length discrepancy and limp
    \item Complications of surgery, including infection and stiffness
    \item Late diagnosis with poorer outcomes
\end{itemize}

\section*{Prognosis and Outcomes}
Early detection transforms the outlook for hip dysplasia. Babies treated with a harness in the first months of life have over 90\% success rates, and normal hip development is usually achieved. Treatment after walking age is more involved but still generally successful, with the majority of children growing up without significant disability. Untreated dysplasia leads to premature arthritis, but adults with symptoms can be helped by physiotherapy, preservation surgery, or replacement. Screening is the single most important factor in good outcomes.

\section*{Prevention and Self-Care}
\begin{itemize}
    \item Attend all newborn and child health checks so hips are examined
    \item Ensure at-risk babies (breech, family history, girls) have ultrasound imaging
    \item Follow the treatment plan if a harness or brace is prescribed
    \item Swaddle babies with the legs free to move
    \item Report limping, hip pain, or leg length differences in older children
    \item Adults with hip pain and a history of dysplasia should seek specialist assessment
    \item Maintain a healthy weight and strong muscles to protect the hips
\end{itemize}

\section*{Sources and Further Reading}
The following reputable sources provide further information for healthcare students and patients seeking reliable, current advice about hip dysplasia.

\begin{itemize}
    \item Healthdirect. \textit{Hip dysplasia in babies.} https://www.healthdirect.gov.au/hip-dysplasia
    \item Raising Children Network. \textit{Developmental dysplasia of the hip.} https://raisingchildren.net.au/
    \item International Hip Dysplasia Institute. \textit{Hip dysplasia information.} https://hipdysplasia.org/
    \item NHS. \textit{Developmental dysplasia of the hip.} https://www.nhs.uk/conditions/developmental-dysplasia-of-the-hip/
    \item Mayo Clinic. \textit{Hip dysplasia: symptoms and causes.} https://www.mayoclinic.org/
    \item Australian Orthopaedic Association. \textit{Hip dysplasia information.} https://aoa.org.au/
\end{itemize}
